\documentclass[../../main.tex]{subfiles}
\begin{document}
\section{向量空间}\label{section:appB-vector-spaces}

设 $\mathbb R$ 表示实数域。$\mathbb R$ 上的向量空间（或实向量空间）是一个集合 $V$，赋予两个运算：向量加法 $V\times V\to V$，记作 $(v,w)\mapsto v+w$；数乘 $\mathbb R\times V\to V$，记作 $(a,v)\mapsto av$，满足下列性质：
\begin{enumerate}[(i)]
\item $V$ 在向量加法下是 Abel 群。
\item 数乘满足下列恒等式：
\begin{align*}
a(bv)=(ab)v\qquad\text{对所有 }v\in V\text{ 与 }a,b\in\mathbb R;
\end{align*}
\begin{align*}
1v=v\qquad\text{对所有 }v\in V.
\end{align*}
\item 数乘与向量加法由下列分配律联系：
\begin{align*}
(a+b)v=av+bv\qquad\text{对所有 }v\in V\text{ 与 }a,b\in\mathbb R;
\end{align*}
\begin{align*}
a(v+w)=av+aw\qquad\text{对所有 }v,w\in V\text{ 与 }a\in\mathbb R.
\end{align*}
\end{enumerate}

这个定义可在两个方向上推广。首先，处处把 $\mathbb R$ 换成任意域 $\mathbb F$，我们得到 $\mathbb F$ 上向量空间的定义。特别地，我们有时有机会考虑 $\mathbb C$ 上的向量空间，称为复向量空间。除非另行说明，所有向量空间都假定是实的。

其次，如果把 $\mathbb R$ 换成交换环 $R$，这就成为 $R$ 上模（或 $R$-模）的定义。例如，若 $\mathbb Z$ 表示整数环，容易验证 $\mathbb Z$ 上的模就是加法下的 Abel 群。

向量空间的元素通常称为\emph{向量}。当需要把它们与向量区分时，基域（除非另行说明为 $\mathbb R$）的元素称为\emph{标量}。

设 $V$ 是向量空间。$V$ 的子集 $W\subset V$ 如果在向量加法与数乘下封闭，它本身在相同运算下就是一个向量空间，称为 $V$ 的一个\emph{子空间}。为避免与拓扑中“子空间”一词的用法混淆，我们有时对这种意义下向量空间的子空间使用\emph{线性子空间}一词，而对赋予子空间拓扑的拓扑空间子集使用\emph{拓扑子空间}一词。

形如 $\sum_{i=1}^k a_i v_i$（其中 $a_i$ 是标量，$v_i\in V$）的有限和称为向量 $v_1,\dots,v_k$ 的一个\emph{线性组合}。（我们把系数 $a_i$ 写上标而非下标的原因，是为了与第 1 章解释的 Einstein 求和约定一致。）如果 $S$ 是 $V$ 的任意子集，$S$ 中元素的所有线性组合所成集合称为 $S$ 的\emph{张成}，记作 $\operatorname{span}(S)$；容易看出它是含 $S$ 的最小子空间。如果 $V=\operatorname{span}(S)$，就说 $S$ \emph{张成} $V$。按约定，不含任何元素的线性组合视为和为零，空集的张成为 $\{0\}$。

如果 $p,q$ 是 $V$ 的点，从 $p$ 到 $q$ 的\emph{线段}是集合 $\{(1-t)p+tq:0\leqslant t\leqslant1\}$。如果对每对点 $p,q\in B$，从 $p$ 到 $q$ 的线段都含于 $B$，则称子集 $B\subset V$ \emph{凸}。

\subsection*{基与维数}

设 $V$ 是向量空间。如果存在形如 $\sum_{i=1}^k a_i v_i=0$ 的线性关系，其中 $v_1,\dots,v_k$ 是 $S$ 的不同元素且至少有一个系数 $a_i$ 非零，则称子集 $S\subset V$ \emph{线性相关}；否则称 $S$ \emph{线性无关}。换言之，$S$ 线性无关，当且仅当 $S$ 中不同元素的唯一和为零的线性组合是所有标量系数都为零的那一个。注意含零向量的每个集合都线性相关。按约定，空集视为线性无关。

经常重要的是处理 $V$ 中向量的有序 $k$ 元组；这样的 $k$ 元组记作 $(v_1,\dots,v_k)$ 或 $(v_i)$，用圆括号而非花括号，以与（无序的）元素集合 $\{v_1,\dots,v_k\}$ 区分。当我们考虑有序 $k$ 元组时，线性相关的含义略有不同。我们说 $(v_1,\dots,v_k)$ 是线性相关的 $k$ 元组，如果存在不全为零的标量 $(a_1,\dots,a_k)$ 使 $\sum_{i=1}^k a_i v_i=0$；否则它是线性无关的 $k$ 元组。线性无关集合与线性无关 $k$ 元组之间的唯一区别是后者不能有重复向量。例如，若 $v\in V$ 是非零向量，则有序对 $(v,v)$ 线性相关，而集合 $\{v,v\}=\{v\}$ 线性无关。另一方面，若 $(v_1,\dots,v_k)$ 是任意线性无关 $k$ 元组，则集合 $\{v_1,\dots,v_k\}$ 也线性无关。

\begin{exercise}\label{exercise:appB-B-1}
设 $V$ 是向量空间。证明下列命题。
\begin{enumerate}[(a)]
\item 如果 $S\subset V$ 线性无关，则 $S$ 的每个子集线性无关。
\item 如果 $S\subset V$ 线性相关或张成 $V$，则 $V$ 中每个真含 $S$ 的子集线性相关。
\item 含多于一个元素的子集 $S\subset V$ 线性相关，当且仅当某个元素 $v\in S$ 可表为 $S\setminus\{v\}$ 中元素的线性组合。
\item 如果 $(v_1,\dots,v_k)$ 是 $V$ 中线性相关 $k$ 元组且 $v_1\neq0$，则某个 $v_i$ 可表为前面向量 $(v_1,\dots,v_{i-1})$ 的线性组合。
\end{enumerate}
\end{exercise}

$V$ 的一个\emph{基}是线性无关且张成 $V$ 的子集 $S\subset V$。如果 $S$ 是 $V$ 的基，$V$ 的每个元素都有唯一方式表为 $S$ 中元素的线性组合。如果 $V$ 有有限基，则称 $V$ \emph{有限维}，否则称\emph{无限维}。平凡向量空间 $\{0\}$ 是有限维的，因为它以空集为基。

如果 $V$ 有限维，$V$ 的一个\emph{有序基}是赋予基向量特定次序的基，等价地说是张成 $V$ 的线性无关 $n$ 元组 $(E_i)$。对大多数目的而言，有序基比无序基更有用，故我们总假定（常不加说明）每个基都带有给定次序。

如果 $(E_1,\dots,E_n)$ 是 $V$ 的（有序）基，每个向量 $v\in V$ 都有唯一方式表为基向量的线性组合：
\begin{align*}
v=\sum_{i=1}^n v^i E_i.
\end{align*}
数 $v^i$ 称为 $v$ 关于此基的\emph{分量}，有序 $n$ 元组 $(v^1,\dots,v^n)$ 称为它的\emph{基表示}。（这是一个需要有序基的定义的例子。）

\begin{lemma}\label{lemma:appB-spanning-dependence}
设 $V$ 是向量空间。如果 $V$ 由 $n$ 个向量的集合张成，则 $V$ 中含多于 $n$ 个向量的每个子集线性相关。
\end{lemma}

\begin{proof}
假设 $\{v_1,\dots,v_n\}$ 是张成 $V$ 的 $n$ 元集合。为证引理，只需证每个恰含 $n+1$ 个向量的集合线性相关。设 $\{w_1,\dots,w_{n+1}\}$ 是这样一个集合。若任一 $w_i$ 为零，则该集合显然相关，故不妨设它们全非零。由习题 B.1(b)，集合 $\{w_1,v_1,\dots,v_n\}$ 线性相关，因而有序 $(n+1)$ 元组 $(w_1,v_1,\dots,v_n)$ 也线性相关。由习题 B.1(d)，其中某个向量 $v_j$ 可表为 $\{w_1,v_1,\dots,v_{j-1}\}$ 的线性组合，故集合 $\{w_1,v_1,\dots,v_{j-1},v_{j+1},\dots,v_n\}$ 仍张成 $V$。必要时重新编号 $v_i$，可设集合 $\{w_1,v_2,\dots,v_n\}$ 张成 $V$。

现归纳假设 $\{w_1,w_2,\dots,w_{k-1},v_k,\dots,v_n\}$ 张成 $V$。如前，$(n+1)$ 元组 $(w_1,w_2,\dots,w_{k-1},w_k,v_k,\dots,v_n)$ 线性相关，故该表中某个向量可表为前面向量的线性组合。如果某个 $w_i$ 可如此写出，则集合 $\{w_1,\dots,w_{n+1}\}$ 相关，我们就完成。否则，某个 $v_j$ 可如此写出，重新排序后可设 $\{w_1,w_2,\dots,w_k,v_{k+1},\dots,v_n\}$ 仍张成 $V$。继续归纳，当到 $k=n$ 时，若尚未表明 $w_i$ 相关，则得出集合 $\{w_1,\dots,w_n\}$ 张成 $V$。但这由习题 B.1(b) 意味着集合 $\{w_1,\dots,w_{n+1}\}$ 线性相关。
\end{proof}

\begin{proposition}\label{proposition:appB-basis-cardinality}
如果 $V$ 是有限维向量空间，则 $V$ 的所有基含相同数目的元素。
\end{proposition}

\begin{proof}
如果 $\{E_1,\dots,E_n\}$ 是 $V$ 的含 $n$ 个元素的基，则引理 B.2 蕴含每个含多于 $n$ 个元素的集合线性相关，故没有基能含多于 $n$ 个元素。另一方面，如果存在含少于 $n$ 个元素的基，则引理 B.2 蕴含 $\{E_1,\dots,E_n\}$ 线性相关，矛盾。
\end{proof}

由于前一命题，如果 $V$ 是有限维向量空间，定义 $V$ 的\emph{维数}（记作 $\dim V$）为任意基中元素的数目是有意义的。

\begin{exercise}\label{exercise:appB-B-4}
假设 $V$ 是有限维向量空间。
\begin{enumerate}[(a)]
\item 证明每个张成 $V$ 的集合都含一个基，$V$ 的每个线性无关子集都含于某个基。
\item 证明每个子空间 $S\subset V$ 有限维且满足 $\dim S\leqslant\dim V$，等号当且仅当 $S=V$。
\item 证明 $\dim V=0$ 当且仅当 $V=\{0\}$。
\end{enumerate}
\end{exercise}

\begin{exercise}\label{exercise:appB-B-5}
假设 $V$ 是无限维向量空间。
\begin{enumerate}[(a)]
\item 用 Zorn 引理证明 $V$ 的每个线性无关子集都含于某个基。
\item 证明 $V$ 的任意两个基有相同基数。[提示：假设 $S$ 与 $T$ 是使 $S$ 基数大于 $T$ 的基。$T$ 的每个元素可表为 $S$ 中元素的线性组合，假设保证 $S$ 中某个元素不出现在 $T$ 元素的任何表达式中。证明该元素可表为 $S$ 中其他元素的线性组合，与 $S$ 线性无关矛盾。]
\end{enumerate}
\end{exercise}

如果 $S$ 是有限维向量空间 $V$ 的子空间，我们定义 $S$ 在 $V$ 中的\emph{余维数}为 $\dim V-\dim S$。由习题 B.4(b)，$S$ 的余维数总非负，且为零当且仅当 $S=V$。一个（线性）\emph{超平面}是余维数为 $1$ 的线性子空间。

\begin{example}[Euclid 空间]\label{example:appB-euclidean-spaces}
对每个整数 $n\geqslant0$，$\mathbb R^n$ 在通常的向量加法与数乘运算下是实向量空间：
\begin{align*}
(x^1,\dots,x^n)+(y^1,\dots,y^n)=(x^1+y^1,\dots,x^n+y^n),
\end{align*}
\begin{align*}
a(x^1,\dots,x^n)=(ax^1,\dots,ax^n).
\end{align*}
$\mathbb R^n$ 有一个自然基 $(e_1,\dots,e_n)$，称为标准基，其中 $e_i=(0,\dots,1,\dots,0)$ 是第 $i$ 位为 $1$、其余为零的向量；因而 $\mathbb R^n$ 有维数 $n$，正如所料。任意元素 $x\in\mathbb R^n$ 可写 $(x^1,\dots,x^n)=\sum_{i=1}^n x^i e_i$，故它关于标准基的分量就是其坐标 $(x^1,\dots,x^n)$。
\end{example}

\begin{example}[复 Euclid 空间]\label{example:appB-complex-euclidean-spaces}
数乘与向量加法恰如实情形定义，$n$ 维复 Euclid 空间 $\mathbb C^n$ 成为复向量空间。由于如上定义的向量 $(e_1,\dots,e_n)$ 构成 $\mathbb C^n$ 在 $\mathbb C$ 上的基，故 $\mathbb C^n$ 作为复向量空间维数为 $n$。

通过把数乘限制为实标量，我们也可把 $\mathbb C^n$ 看作实向量空间。在此情形中，容易验证向量 $(e_1,\mathrm{i}e_1,\dots,e_n,\mathrm{i}e_n)$ 构成 $\mathbb C^n$ 在 $\mathbb R$ 上的基，故 $\mathbb C^n$ 作为实向量空间维数为 $2n$。
\end{example}

如果 $S$ 与 $T$ 是向量空间 $V$ 的子空间，记号 $S+T$ 表示所有形如 $v+w$（其中 $v\in S$、$w\in T$）的向量所成集合。容易看出它是 $V$ 的子空间，实际上是 $S\cup T$ 张成的子空间。如果 $S+T=V$ 且 $S\cap T=\{0\}$，则称 $V$ 是 $S$ 与 $T$ 的（内）直和，写作 $V=S\oplus T$。两个线性子空间 $S,T\subset V$ 如果 $V=S\oplus T$，就称为互补子空间。在此情形中，$V$ 中每个向量都有唯一方式表为 $S$ 的元素加 $T$ 的元素。

\begin{exercise}\label{exercise:appB-B-8}
假设 $S$ 与 $T$ 是有限维向量空间 $V$ 的子空间。
\begin{enumerate}[(a)]
\item 证明 $S\cap T$ 是 $V$ 的子空间。
\item 证明 $\dim(S+T)=\dim S+\dim T-\dim(S\cap T)$。
\item 假设 $V=S+T$。证明 $V=S\oplus T$ 当且仅当 $\dim V=\dim S+\dim T$。
\end{enumerate}
\end{exercise}

\begin{exercise}\label{exercise:appB-B-9}
设 $V$ 是有限维向量空间。证明每个子空间 $S\subset V$ 在 $V$ 中都有互补子空间。事实上，给定 $V$ 的任意基 $(E_1,\dots,E_n)$，证明存在整数 $\{1,\dots,n\}$ 的某个子集 $\{i_1,\dots,i_k\}$，使 $\operatorname{span}(E_{i_1},\dots,E_{i_k})$ 是 $S$ 的补。[提示：选 $S$ 的基 $(F_1,\dots,F_m)$，对有序 $(m+n)$ 元组 $(F_1,\dots,F_m,E_1,\dots,E_n)$ 应用习题 B.1(d)。]
\end{exercise}

假设 $S\subset V$ 是线性子空间。$V$ 中形如 $v+S=\{v+w:w\in S\}$（对某个固定 $v\in V$）的任意子集称为 $V$ 中平行于 $S$ 的\emph{仿射子空间}。如果 $S$ 是线性超平面，则平行于 $S$ 的任意仿射子空间称为\emph{仿射超平面}。

\begin{exercise}\label{exercise:appB-B-10}
设 $V$ 是向量空间，$v+S$ 是 $V$ 中平行于 $S$ 的仿射子空间。
\begin{enumerate}[(a)]
\item 证明 $v+S$ 是线性子空间，当且仅当它含 $0$，而这当且仅当 $v\in S$。
\item 证明 $v+S=\widetilde v+\widetilde S$，当且仅当 $S=\widetilde S$ 且 $v-\widetilde v\in S$。
\end{enumerate}
\end{exercise}

由于前一习题的 (b)，我们可无歧义地把 $v+S$ 的维数定义为 $S$ 的维数。

对每个向量 $v\in V$，仿射子空间 $v+S$ 也称为由 $v$ 确定的 $S$ 的\emph{陪集}。$S$ 的所有陪集所成集合 $V/S$ 称为 $V$ 模 $S$ 的\emph{商}。

\begin{exercise}\label{exercise:appB-B-11}
假设 $V$ 是向量空间，$S$ 是 $V$ 的线性子空间。用
\begin{align*}
(v+S)+(w+S)=(v+w)+S,\qquad c(v+S)=(cv)+S
\end{align*}
定义陪集的向量加法与数乘。
\begin{enumerate}[(a)]
\item 证明商 $V/S$ 在这些运算下是向量空间。
\item 证明如果 $V$ 有限维，则 $\dim V/S=\dim V-\dim S$。
\end{enumerate}
\end{exercise}
\end{document}
